\section{Conclusion}
\label{sec:conclusion}

We have shown that LLMs create unexpected partial isomorphisms---structure-preserving mappings between concept domains that differ from human intuitions.
Through embedding analysis, we identified 162 significantly surprising cross-domain concept pairs.
Through behavioral probing, we demonstrated that these unexpected analogies improve reasoning quality.
Through mechanistic analysis, we provided evidence that surprising pairs share activation patterns concentrated in semantic processing layers.

The key takeaway is that neural compression does not merely entangle features for efficiency---it discovers genuine structural parallels across domains that humans have not commonly articulated.
These ``alien analogies'' reflect functional and structural properties of the underlying concepts, suggesting that large-scale text compression creates emergent conceptual organization.

Future work should validate these findings with human similarity judgments, scale the activation analysis across model sizes to study how partial isomorphisms evolve with capacity, and use causal interventions (activation patching) to verify that shared representations are functionally important rather than merely correlational.
